\section{Experiments and Recorded Measurements}
\label{sec:experiments}

\subsection{GEN-003: Current Four-Model Throughput Record}
\recordstatus{Measured; supplied summary, incomplete provenance.}
\evidence{Suite \texttt{comparison-v2}; received 2026-10-22. Execution date
not supplied. Tables are generated from the frozen JSON transcription.}

\textbf{Question.} What throughput values were recorded for the four model
variants on the selected in-distribution and out-of-distribution categories?

\textbf{Setup.} The source is the suite owner's supplied
\texttt{summary.md}. Tables~\ref{tab:nemotron} and~\ref{tab:redhat} preserve all
28 reported measurements without adding older model rows. Step counts follow
the supplied labels. Correctness was not checked. This record does not include
the native per-request outputs, latency statistics, acceptance counters, or
server provenance for this suite.

\input{generated/throughput_tables}

\textbf{Record notes.} The values are recorded separately by category and
checkpoint. We do not pool categories into an overall score or infer a final
model ranking. The checkpoint pair provides observations at two reported
training stages, not a repeated or controlled causal estimate of additional
training. The two tables do not establish output quality or a target-only
decoding speedup.

\textbf{Missing evidence.} Native GuideLLM results, effective configurations,
hardware metadata, prompt fingerprints, checkpoint hashes, and repeated
measurements remain to be attached to this suite record.

\subsection{SYS-001: Cold and Warm Activation Reuse}
\recordstatus{Measured; console timings transcribed in the cache record.}
\evidence{Pilot logs received 2026-10-16; reported A100 80 GB, 400 validation
payloads evaluated in context-4 groups.}

\textbf{Question.} What stage wall times were observed before and after
activation payloads were available in the shared RAM cache?

\input{generated/cache_table}

\textbf{Record notes.} The payloads occupied 16.81 GiB. Subsequent jobs in the
same process reported initial validation passes of 19.99 and 16.86 seconds.
The first pass includes extraction and startup effects, whereas the warm pass
still performs generator/drafter evaluation. The difference is an end-to-end
observation, not an isolated cache microbenchmark or inference speedup. In a
direct-training pilot the cache reached 63.97 GiB, with 365 evictions and
69.67 GiB final process RSS. Later prefetch and summary-cache changes have no
separate GPU timing comparison in the evidence included here.

\textbf{Missing evidence.} Repeated cold/warm trials with separated extraction,
construction, transfer, and compute timings remain available follow-up work.

\subsection{GEN-002: Output Repeatability Investigation}
\recordstatus{Investigation deferred; exact sequence comparisons recorded.}
\evidence{Local correctness reports summarized in the Markdown registry,
observation dated 2026-10-21. Raw reports are not archived with this report.}

\textbf{Question.} Were target-only and speculative outputs repeatable under
the supplied prompt/settings fingerprints?

On 50 Nemotron math prompts, original versus repeated target-only decoding
matched exactly for 23/50 sequences; original versus repeated EAGLE3 decoding
matched for 50/50. Repeated target-only versus repeated EAGLE3 matched for
17/50. These are whole-sequence match counts, not token-error rates.

\textbf{Record notes.} The cause remains unresolved. Runtime environment,
numerical behavior, and decoding implementation differences have not been
isolated. Current timing experiments explicitly record correctness as unchecked
or unverified. No correctness claim is inferred from the throughput tables.

\textbf{Missing evidence.} Preserve and compare complete runtime provenance and
repeat a controlled reference experiment before making an exact-output claim.

\subsection{BATCH-001: Shared-Adapter Concurrency Sweep}
\recordstatus{Implemented; GPU measurements pending.}
\evidence{Implementation revision \texttt{c0ee0e4}; CPU lifecycle and HTTP
tests, plus an audit of the pinned vLLM hook upgrade.}

\textbf{Question.} What measurements arise when a generator conditions one
shared adapter on a fixed group at concurrency 1, 2, 4, or 8?

\textbf{Planned measurement configuration.} Reuse the frozen category prompts
for all model/concurrency variants. Record generated output tokens/s,
speculative acceptance where available, request latency, actual mean batch size,
and generator/merge/restoration costs. Keep the adapter fixed while a group
drains. Original-drafter and static-LoRA variants use matching closed groups.

\textbf{Available checks.} Tests cover packed prompt boundaries, one generator
invocation per group, restoration after the last member, delayed HTTP cleanup,
closed admission, partial groups, and suite resume across concurrency variants.
The initial broader CPU/HTTP run passed 94 tests; a later focused run after
controller refinements passed 50. These runs overlap and are not additive.

\textbf{Pending evidence.} No A100 execution, batched throughput table,
acceptance comparison, or memory measurement is supplied. CPU checks establish
neither GPU compatibility under load nor a performance benefit.
